\documentclass[../main]{subfiles}

\begin{document}

\chapter{Primer Principio de la Termodinámica}

\section{Introducción al Primer Principio de la Termodinámica}

El Primer Principio de la Termodinámica es una ley fundamental en la física que establece la conservación de la energía en sistemas termodinámicos. Esta ley establece que la energía total de un sistema cerrado es constante, lo que significa que la energía no puede ser creada ni destruida, solo transformada de una forma a otra.

\info{Primer Principio de la Termodinámica}{
El Primer Principio de la Termodinámica establece que la energía interna de un sistema termodinámico es igual a la suma del calor agregado al sistema y el trabajo realizado por el sistema sobre su entorno.
\begin{equation}
    \Delta U = Q + W
\end{equation}

Donde:
\begin{itemize}
    \item $\Delta U$: Cambio en la energía interna del sistema [J].
    \item $Q$: Calor agregado al sistema [J].
    \item $W$: Trabajo realizado por el sistema sobre su entorno [J].
\end{itemize}
}

\actividad{Responder el siguiente cuestionario}
\begin{enumerate}
    \item ¿Qué establece el Primer Principio de la Termodinámica?
    \item ¿Cómo se relaciona la energía interna de un sistema con el calor y el trabajo?
    \item ¿Qué significa que la energía sea una magnitud conservativa?
    \item ¿Qué condiciones deben cumplirse para que se aplique el Primer Principio de la Termodinámica?
    \item ¿Cuál es la importancia del Primer Principio de la Termodinámica en el estudio de los sistemas físicos?
    \item ¿Cuál es la diferencia entre un sistema cerrado y un sistema abierto en el contexto del Primer Principio de la Termodinámica?
    \item ¿Qué es el calor sensible y cómo se relaciona con el cambio de temperatura de un sistema?
    \item ¿Qué es el calor latente y cuál es su importancia en procesos termodinámicos?
    \item ¿Cómo se relaciona el Primer Principio de la Termodinámica con la ley de conservación de la energía?
    \item ¿Qué ejemplos cotidianos pueden ilustrar la aplicación del Primer Principio de la Termodinámica?
\end{enumerate}

\actividad{Resolver los siguientes problemas}
\begin{enumerate}
    \item Un sistema absorbe $500 \, \text{J}$ de calor y realiza un trabajo de $200 \, \text{J}$. ¿Cuál es el cambio en su energía interna?
    \item Un gas ideal a $300 \, \text{K}$ y $1 \, \text{atm}$ ocupa un volumen de $0.1 \, \text{m}^3$. Si se duplica la presión a temperatura constante, ¿cuál será su nuevo volumen?
    \item Un sistema realiza un trabajo de $150 \, \text{J}$ y libera $250 \, \text{J}$ de calor al entorno. ¿Cuál es el cambio en su energía interna?
    \item En un proceso isocórico, se añaden $400 \, \text{J}$ de calor a un gas ideal. ¿Cuál es el cambio en su energía interna?
    \item Un gas ideal realiza un trabajo de $100 \, \text{J}$ cuando su volumen cambia de $0.5 \, \text{m}^3$ a $1 \, \text{m}^3$ a presión constante. ¿Cuál es su presión?
    \item Un sistema recibe $600 \, \text{J}$ de calor y su energía interna aumenta en $400 \, \text{J}$. ¿Cuánto trabajo realiza el sistema?
    \item Un gas ideal a $273 \, \text{K}$ ocupa un volumen de $22.4 \, \text{L}$. Si se duplica la temperatura a presión constante, ¿cuál será su nuevo volumen?
    \item Un sistema tiene un cambio en su energía interna de $-300 \, \text{J}$ y realiza un trabajo de $100 \, \text{J}$. ¿Cuánto calor ha intercambiado con su entorno?
    \item Un gas ideal a $500 \, \text{K}$ y $2 \, \text{atm}$ ocupa un volumen de $0.05 \, \text{m}^3$. Si se reduce la temperatura a la mitad a presión constante, ¿cuál será su nuevo volumen?
    \item Un sistema absorbe $800 \, \text{J}$ de calor y realiza un trabajo de $300 \, \text{J}$. Si su cambio en la energía interna es de $200 \, \text{J}$, ¿cuál es el signo del calor transferido?
\end{enumerate}

\end{document}
